\chapter{Human-Building Interaction Meets Affective Experience}\label{ch:background}

\chapterepigraph{A building is not something you finish. A building is something you start.}{Stewart Brand}

\begin{synopsis}
This chapter positions the dissertation across Human-Computer Interaction, Human-Building Interaction, affective interaction, biophilic design, and architectural experience.
\end{synopsis}

\section{human-building interaction}
\marginterm{HBI}{Human-Building Interaction studies the relation between people, buildings, embedded technologies, and organisational/infrastructural systems.} Human-Building Interaction extends HCI into built environments: buildings are not just contexts for interaction, but interactional systems themselves.

\section{affective interaction}
Affective interaction shifts attention from detecting emotional labels to understanding how affect is enacted, interpreted, sustained, and negotiated in practice.\sidenote{This is a good place for a compact theoretical clarification or a citation.}

\section{biophilic technologies}
Biophilic design is often treated as a matter of plants, daylight, water, and natural materials. A biophilic technology perspective asks how computational systems might mediate experiences of nature, atmosphere, care, restoration, and environmental relation.

\begin{figure}
\begin{fullwidth}
\centering
\fbox{\parbox{0.86\linewidth}{\centering\vspace{1.5in}Full-width conceptual diagram placeholder\vspace{1.5in}}}
\caption{Use \texttt{fullwidth} for large conceptual diagrams that need the full page width.}
\end{fullwidth}
\end{figure}
